\scanpage{5}
of the Artin ring $B_n=H^o(X_n,\underline{O}_{X_n})$. Assuming $U$ hence $X$ hence $X_n$ to be connected, this Artin ring is local, and it's Brauer group coincides with the Brauer group of its residue field $k_n$. Of course, \deleted{thus} for an element of $P(k)=\lim P_n(k)$, these vari\deleted{\unsure{e}}ous obstructions match together and come from an element of the Brauer group of the residue fi\rlap{l}e\rlap{e}ld $k'=\handwritten{\bigcap_n}k_n$ of the ring $H^o(X,\underline{O}_X)=\underset{\handwritten{n}}{\underset{\handwritten{\longleftarrow}}{\lim}}H^o(X_n,\underline{O}_{X_n})$, \deleted{of} which is just the normalisation of $A$. If for instance $A$ is \uline{normal}, \tikz[baseline=(s.base)]{\node[inner sep=1pt](s){$k'=k$};\node[anchor=south west,inner sep=1pt](a) at ([xshift=1pt,yshift=7pt]s.north west){\inserted{\handwritten{then}}};\draw ([xshift=-5pt,yshift=-3pt]s.south west) to[out=90,in=190] ([xshift=1pt,yshift=-2pt]a.south west)--([xshift=1pt,yshift=-2pt]a.south east) to[out=0,in=-90] ([xshift=3pt,yshift=4pt]a.south east);} and the obstruction lies in $\operatorname{Br}(k)$.

Whatever this be, defining $\unsure{Q|\mathbb{Q}}$ as before, we still do get an injective homomorphism (6), but even if (5) is injective, we cannot be sure that (6) is bijective. We are however if $k$ \unsure{is|ss} separable closed. Therefore, we still \deleted{get} can more or less describe $\mathord{\rlap{$\deleted{P}$}\handwritten{Q}}$ (or at least $\mathord{\rlap{$\deleted{P}$}\handwritten{Q}}(k)$) in termes of local Picard groups, in terms of \tikz[baseline=(s.base)]{\node[inner sep=1pt](s){descent};\node[anchor=south west,inner sep=1pt](a) at ([xshift=1pt,yshift=7pt]s.north west){\inserted{Galois}};\draw ([xshift=-5pt,yshift=-3pt]s.south west) to[out=90,in=190] ([xshift=1pt,yshift=-2pt]a.south west)--([xshift=1pt,yshift=-2pt]a.south east) to[out=0,in=-90] ([xshift=3pt,yshift=4pt]a.south east);} from the completion $A'$ of the strict henselization of $A$: if $U'=S'-\handwritten{\{}s'\handwritten{\}}$, where $S'=\operatorname{Spec}(A\mathord{\rlap{${}'$}(})$, \tikz[baseline=(s.base)]{\node[inner sep=0pt](s){and};\fill ([xshift=0pt,yshift=3pt]s.north) circle (.35pt);} if
\[
\begin{aligned}
(9)\qquad\qquad & \handwritten{\pi}=\operatorname{Gal}(\bar{k}/k)\ ,
\end{aligned}
\]
we will have an isomorphism
\[
\begin{aligned}
(10)\qquad\qquad & \deleted{\operatorname{Pic}}\unsure{\mathord{\rlap{$P$}Q}|\mathbb{Q}}(k)\handwritten{\simeq}\underline{\operatorname{Pic}}(U')^{\handwritten{\pi}}\ ,
\end{aligned}
\]
\deleted{This is perf provided $D'$ is torsion free, for instance (5)\,\tikz[baseline=-.3ex,x=1em,y=1em]{\fill (.12,-.35) circle (.35pt);\fill (.23,-.20) circle (.35pt);} injective\unsure{,|.}}

This is perfectly natural and in accordance with the familiar phenomena when dealing with \uline{global} Picard schemes. Therefore $Q$ looks like \rlap{a}the natural object to be called a \deleted{Picard} local Picard scheme.

To make sure it makes a sense, one will have however to analyze somewhat the inverse system $(P_n)$, and the morphism (5). Moreover, to feel really secure, one should check that up to canonical isomorphism, the group scheme $Q$ together with (10) does not depend on the choice of
% archive sheet number: 4
